\begin{tikzpicture}[font=\small,>=stealth,line join=round]
\draw[draw=teal!45!black,fill=green!3,rounded corners=3pt] (0,0) rectangle (13.6,6.1);
\node at (6.8,5.55) {$a-t$ 中的负号：$t$ 增大时，这个输入向左移动};
\begin{scope}[shift={(1.3,1.6)},x=1.2cm,y=0.65cm]
\draw[gray,->] (-0.1,0)--(4.4,0) node[right] {$x$};
\draw[gray,->] (0,-0.2)--(0,4.8);
\draw[teal!75!black,very thick,domain=0.3:3.7,samples=100] plot(\x,{(\x-2)^2+1});
\fill[teal!75!black] (1,2) circle(2.5pt);
\fill[teal!75!black] (3,2) circle(2.5pt);
\draw[dashed,gray] (2,0)--(2,4.5);
\draw[->,orange!85!black,very thick] (1,-.8)--(.55,-.8);
\draw[->,blue!70!black,very thick] (3,-.8)--(3.45,-.8);
\node[below] at (1,0) {$a-t=1$};\node[below] at (3,0) {$a+t=3$};
\end{scope}
\node[anchor=west,align=left,text width=5.4cm] at (7.3,3) {取 $a=2,t=1$。\\[10pt]$t$ 增大时：\\左点向左，右点向右；\\两边都沿图像升高。\\[10pt]但 $f^{\prime}(1)=-2$、$f^{\prime}(3)=2$。\\负号来自左边输入的反向变化。};
\node at (6.8,.4) {橙箭头：$a-t$ 减小；蓝箭头：$a+t$ 增大。};
\end{tikzpicture}
